% US Dollar, October 2026 outlook: Bristol Trading Society Currency Report Card
% Compile with pdflatex (twice is not needed). Edit the text and numbers below;
% the look lives in bts-reportcard.sty. Accent options: red (default), olive, ink.
\documentclass[10pt]{article}
\usepackage[red]{bts-reportcard}

\currency{US Dollar}
\subtitle{October 2026 outlook: firm into the FOMC, fraying beyond it}
\analyst{Valeriy Sinko}
\asof{28 Sep 2026}
\issue{1}
\tacticalview{USD bullish}
\structuralview{Neutral}
\logo{bts-logo}

\begin{document}
\makereportheader

%% Row 1: 1 to 3 month outlook | indicators panel ------------------------------
\reportrow{%
  \begin{outlook}{1 to 3 month outlook}{The market underprices a second Fed hike; the dollar stays firm into the 28 Oct FOMC.}
    \point{Priced}{72\% odds of a 25bp hike on 28 Oct, but only $\sim$30\% for a second by December (CME FedWatch, fed funds futures, 25--28 Sep). Fed median: 4.1\% by year end, one more hike then a pause.}
    \point{View}{\textbf{Bullish USD:} the Fed hikes more than priced.
      \evidence{Fed}{16 of 18 officials expect another hike, four see two; the 2026 median rose from 3.8\% to 4.1\% (SEP, 16 Sep).}
      \evidence{Inflation}{CPI 3.4\%, core +0.3\% m/m (BLS, Aug); Brent near \$107; consumers expect 4.6\% (UMich, Sep).}
      \evidence{Jobs}{payrolls +162k vs +53k expected; unemployment 4.1\% (BLS, Aug): the economy can bear higher rates.}}
    \point{Positioning}{CFTC DXY net longs +10,330 (22 Sep), near a 3-month low: not crowded, room for buyers.}
    \point{Change our mind}{Payrolls under 50k (2 Oct), core CPI $\leq$0.2\% m/m (14 Oct), a hold on 28 Oct, or an Iran ceasefire that knocks oil lower.}
  \end{outlook}%
}{%
  \begin{indicators}
    % Comparisons are against the DXY basket: EUR 57.6%, JPY 13.6%, GBP 11.9%,
    % CAD 9.1%, SEK 4.2%, CHF 3.6%. Weighted policy rate = 2.33% (28 Sep 2026).
    \indicator{DXY spot (25 Sep)}{100.97}
    \indicator{DXY 3m forward}{\fxforward[2]{100.97}{3.875}{2.3345}{0.25}}
    \indicator{Fed funds (16 Sep)}{3.75--4.00\%}
    \indicator{Policy gap vs DXY basket}{+1.54 pts}
    \indicator{Policy gap vs EUR (trade)}{+1.375 pts}
    \indicator{Next FOMC \textperiodcentered{} priced}{28 Oct: 72\%}
    \indicator{2y Treasury; 1m chg}{[4.81\%] \textperiodcentered{} +47bp}
    \indicator{Real policy rate (US)}{+0.5\%}
    \indicator{CPI vs 2\% target (Aug)}{3.4\% (+1.4 pts)}
    \indicator{Valuation: REER (May)}{114.0}
    \lastindicator{CFTC DXY net (22 Sep)}{+10,330}
  \end{indicators}%
}

%% Row 2: 6 to 12 month outlook | chart 1 -------------------------------------
\reportrow{%
  \begin{outlook}{6 to 12 month outlook}{The rate premium peaks; valuation and fiscal strain start to bite.}
    \point{Priced}{A plateau near 4.1\% through 2027, easing to 3.9\% in 2028 (Fed SEP).}
    \point{Supportive}{Positive real rates; net energy exporter with Brent near \$107; safe-haven demand on Iran.}
    \point{Against}{REER 114, near its record; $\sim$17\% rich to EUR on PPP; a 5\% 10y and Treasury buybacks leaning on yields.}
    \point{Where we disagree}{The market assumes the Fed can outrun a Treasury capping yields.}
  \end{outlook}
  \vspace{4pt}
  \drivertagpos{Real rates +} \drivertagpos{Terms of trade +} \drivertagpos{Dollar smile +}
  \drivertagneg{Fiscal $-$} \drivertagneg{Valuation $-$}%
}{%
  \begin{chartbox}{Chart 1}{The dollar rallied into the hike, back to July highs}%
    {DXY daily close. Source: Portfolio Terminal, 25 Sep 2026. Swap in a 1--2 year DXY vs 2y spread chart before publishing.}
    \begin{tikzpicture}
      \begin{axis}[btschart,
          symbolic x coords={21 Sep,22,23,24,25 Sep}, xtick=data,
          ymin=100.3, ymax=101.45, ytick={100.5,101.0},
          yticklabel style={/pgf/number format/fixed,/pgf/number format/precision=1,/pgf/number format/fixed zerofill}]
        \addplot[mark=*, mark size=1.8pt] coordinates
          {(21 Sep,100.43) (22,100.60) (23,101.10) (24,101.29) (25 Sep,100.97)};
        \addplot[only marks, mark=*, mark size=3pt, color=accent,
                 nodes near coords, point meta=explicit symbolic,
                 nodes near coords style={above, font=\scriptsize\ttfamily}]
          coordinates {(24,101.29) [101.29]};
      \end{axis}
    \end{tikzpicture}
  \end{chartbox}%
}

%% Row 3: trade idea | chart 2 -------------------------------------------------
\reportrow{%
  \begin{tradeidea}{Sell EUR/USD (long USD), spot}
    \traderow{Entry}{\texttt{1.1400}}
    \traderow{Stop}{\texttt{1.1550}}
    \traderow{Target}{\texttt{1.1100}}
    \traderow{Reward : risk}{\texttt{2.0} before carry \textperiodcentered{} \texttt{2.3} after (+39 pips)}
    \traderow{Size}{1\% of book \textperiodcentered{} $\sim$EUR 667k}
    \traderow{Catalysts}{2 Oct payrolls \textperiodcentered{} 14 Oct CPI \textperiodcentered{} 28 Oct FOMC}
    \lasttraderow{Horizon}{Small, short dated; exit after 28 Oct if flat}
  \end{tradeidea}%
}{%
  \begin{chartbox}{Chart 2}{The Fed plans to stay high for two years, then ease}%
    {Fed funds midpoint now vs median projections. Source: FOMC Summary of Economic Projections, 16 Sep 2026.}
    \begin{tikzpicture}
      \begin{axis}[btschart,
          symbolic x coords={Now,2026,2027,2028,2029,Long run}, xtick=data,
          ymin=3.0, ymax=4.45, ytick={3.5,4.0},
          yticklabel={\pgfmathprintnumber[fixed,precision=1,fixed zerofill]{\tick}\%},
          nodes near coords, point meta=explicit symbolic,
          nodes near coords style={above, font=\scriptsize\ttfamily}]
        \addplot[dashed, line width=0.8pt, mark=*, mark size=2pt, mark options={solid}]
          coordinates {(Now,3.875)[3.875] (2026,4.1)[4.1] (2027,4.1)[4.1]
                       (2028,3.9)[3.9] (2029,3.6)[3.6] (Long run,3.2)[3.2]};
        \addplot[only marks, mark=*, mark size=3pt, color=accent] coordinates {(Now,3.875)};
      \end{axis}
    \end{tikzpicture}
  \end{chartbox}%
}

%% Forecasts against the forward ----------------------------------------------
%% Forwards are calculated live from spot and policy rates with \fxforward.
\begin{forecasts}
  \forecastrow{EUR/USD}{3 months}{1.1403}{\fxforward{1.1403}{2.50}{3.875}{0.25}}{1.1200}{\textbf{USD stronger} than the forward}
  \forecastrow{EUR/USD}{12 months}{1.1403}{\fxforward{1.1403}{2.50}{3.875}{1}}{1.1600}{In line with the forward}
  \forecastrow{USD/JPY}{3 / 12 months}{157.25}{\fxforward[2]{157.25}{3.875}{1.25}{0.25} / \fxforward[2]{157.25}{3.875}{1.25}{1}}{[\,] / [\,]}{From the JPY page}
  \lastforecastrow{USD/CAD}{3 / 12 months}{1.4141}{\fxforward{1.4141}{3.875}{2.25}{0.25} / \fxforward{1.4141}{3.875}{2.25}{1}}{[\,] / [\,]}{From the CAD page}
\end{forecasts}

\sourcesnote{Forwards from spot and policy rates as a proxy (USD 3.875\%, EUR 2.50\%, GBP 3.75\%, CHF 0\%, SEK 1.75\%, JPY 1.25\%, CAD 2.25\%; DXY-weighted basket 2.33\%). Sources: Federal Reserve, ECB, BoE, SNB, Riksbank, BLS, BEA, CME FedWatch, CFTC, BIS/FRED, Portfolio Terminal, CNBC; data 16--28 Sep 2026. Values in [brackets] to be verified or completed. Student research for teaching purposes; not investment advice.}

\end{document}
